% MyHatTileStuff/grow_symmetric2.c — Detailed Technical Explanation
% Compile with: pdflatex -shell-escape grow_symmetric2.tex
%
\documentclass[10pt,a4paper]{article}

%% ---------- Packages ----------
\usepackage[T1]{fontenc}
\usepackage[utf8]{inputenc}
\usepackage{geometry}
\geometry{margin=2cm, top=2.2cm, bottom=2.5cm}

\usepackage{microtype}
\usepackage{parskip}
\usepackage{xcolor}
\usepackage{booktabs}
\usepackage{amsmath,amssymb}
\usepackage{graphicx}
\usepackage{hyperref}
\usepackage{enumitem}
\setlist{noitemsep, topsep=2pt}
\usepackage{needspace}

%% minted — syntax highlighting (requires -shell-escape and pygments)
\usepackage{minted}
\setminted{
  fontsize=\footnotesize,
  baselinestretch=1.1,
  breaklines=true,
  linenos=false,
  frame=none,
  bgcolor=codebg
}
\definecolor{codebg}{HTML}{F5F5F5}

%% tcolorbox — side-by-side description/code panels
\usepackage[most]{tcolorbox}
\tcbuselibrary{minted,skins,breakable}

%% Colours
\definecolor{titlecolor}{HTML}{1A4D8F}
\definecolor{headercolor}{HTML}{2E6DB4}
\definecolor{rulecolor}{HTML}{AECBEF}
\definecolor{seedorange}{HTML}{F39C12}
\definecolor{seedred}{HTML}{E74C3C}
\definecolor{seedpurple}{HTML}{8E44AD}
\definecolor{grownblue}{HTML}{A8D8EA}
\definecolor{centerred}{HTML}{C0392B}

%% Side-by-side box style
\tcbset{
  dualbox/.style={
    enhanced,
    colframe=rulecolor,
    colback=white,
    fonttitle=\bfseries\small,
    breakable,
    boxrule=0.6pt,
    arc=3pt,
    sidebyside,
    sidebyside align=top seam,
    sidebyside gap=6mm,
    lefthand ratio=0.46,
    before skip=8pt,
    after skip=8pt,
  }
}

%% Section colours
\usepackage{titlesec}
\titleformat{\section}{\large\bfseries\color{titlecolor}}{}{0em}{}[\nobreak\vspace{-2pt}\color{rulecolor}\hrule\nobreak\vspace{4pt}]
\titleformat{\subsection}{\normalsize\bfseries\color{headercolor}}{}{0em}{}

%% Keep headings attached to the content that follows them
\newcommand{\subsectionbreak}{\nobreak}
\let\polysection\section
\renewcommand\section{\needspace{0.16\textheight}\polysection}
\let\polysubsection\subsection
\renewcommand\subsection{\needspace{0.12\textheight}\polysubsection}

%% Convenient macros
\newcommand{\code}[1]{\texttt{\small #1}}
\newcommand{\file}[1]{\texttt{\footnotesize #1}}
\newcommand{\vcoord}[3]{$(#1,\,#2,\,#3)$}

%% ---------- Document ----------
\begin{document}

\begin{center}
  {\Huge\bfseries\color{titlecolor} grow\_symmetric2.c}\\[4pt]
  {\large\color{headercolor} Growing 3-Fold Symmetric Hat Clusters from the $v[12]$ Pinwheel}\\[6pt]
  {\small \file{MyHatTileStuff/grow\_symmetric2.c} — Technical Reference}\\[2pt]
  {\small\color{gray} \today}
\end{center}

\vspace{6pt}
\noindent\rule{\linewidth}{1pt}
\vspace{4pt}

% -----------------------------------------------------------------------
\section{Purpose}

\code{grow\_symmetric2} enumerates hat-tile clusters that grow ring-by-ring
around a fixed \emph{3-hat pinwheel} seed: three hat tiles sharing the
single vertex $v[12]$ of the base hat cycle, placed at relative orientations
$0°$, $120°$, and $240°$.  After growing $N$ concentric rings the program
keeps only those completions whose outer boundary is invariant under a $120°$
rotation around the shared vertex, outputting them as a \code{.dat} record
file and (optionally) an SVG grid.

\medskip
\noindent The program answers the question: \emph{can a 3-fold rotationally
symmetric hat cluster centred at vertex $v[12]$ be extended indefinitely, or
does the RL0 vertex-arc rulebook force a dead end?}

% -----------------------------------------------------------------------
\section{Background}

\subsection{The Tetrille Lattice and Hat Tile}

The \textbf{tetrille} lattice tiles the plane with equilateral triangles and
squares, producing vertices of three valences: $v=6$ (hexagonal, shared by
two squares and two triangles), $v=4$ (square, shared by one square and two
triangles), and $v=3$ (triangular, shared by one square and one triangle).
A lattice point is an integer triple $(v, x, y)$ where $v$ carries the
valence tag; arithmetic across valence tags requires an explicit lifting
formula.

The hat tile is a 14-vertex polygon whose boundary cycle in \file{hat.tile}
reads:

\begin{minted}{text}
v[0]  = (6, 0, 0)     v[1]  = (4,-1, 0)     v[2]  = (3,-1, 0)
v[3]  = (4,-1,-1)     v[4]  = (6, 0,-1)     v[5]  = (4, 0,-3)
v[6]  = (3,-1,-2)     v[7]  = (4, 1,-3)     v[8]  = (3, 0,-2)
v[9]  = (4, 1,-2)     v[10] = (6, 1,-1)     v[11] = (4, 2,-1)
v[12] = (3, 1,-1)     v[13] = (4, 1,-1)
\end{minted}

\subsection{Choosing the Pinwheel Centre}

A 3-hat pinwheel requires a vertex $p$ with interior angle $120°$ inside the
hat.  Three of the hat's 14 vertices satisfy this: $v[0]$, $v[6]$, and
$v[12]$.  They differ in two important ways.

\begin{center}
\begin{tabular}{@{}llllp{6cm}@{}}
\toprule
\textbf{Vertex} & \textbf{Coord} & \textbf{Valence} & \textbf{RL0 status}
  & \textbf{Notes} \\
\midrule
$v[0]$ & $(6,0,0)$ & hexagonal & \textbf{valid} &
  Tetrille coordinate origin; fixed point of all lattice rotations;
  no translation needed for the symmetry check.  Dead end at level 3
  (51-tile cluster, no further symmetric extension). \\[4pt]
$v[6]$ & $(3,-1,-2)$ & triangular & \textbf{forbidden} &
  Geometrically valid $120°$ angle, but the RL0 vertex-arc dictionary
  prohibits three hats meeting there in a pinwheel arrangement.
  Returns 0 completions even at level 1. \\[4pt]
$v[12]$ & $(3,1,-1)$ & triangular & \textbf{valid} &
  The centre used by this program.  Not the origin, so a translation
  correction is required after each $60°$-lattice rotation.
  Gives an open cluster: 2, 4, 10 symmetric completions at
  levels 1, 2, 3. \\
\bottomrule
\end{tabular}
\end{center}

% -----------------------------------------------------------------------
\section{Rotation Around a Non-Origin Centre}

\subsection{The Lattice Rotation Primitive}

\code{cycle\_transform\_lattice(src, dst, lattice, t)} applies the
\emph{origin-centred} 60°-step lattice rotation
\[
  R_t: (x, y) \;\longmapsto\; (-y,\; x+y)^t
\]
to every vertex of \code{src}.  The transform is exact on all three valence
systems.  A single application with $t=1$ rotates by $60°$; $t=2$ rotates
by $120°$; $t=4$ rotates by $240°$.  The rotation centre is always
\emph{the lattice origin} $(v, 0, 0)$.

\subsection{Rotating Around an Arbitrary Centre}

To rotate a cycle by $120°$ around a point $C$ that is \emph{not} the
origin, the standard conjugation trick is used:
\[
  \text{Rot}_{C,120°}(\mathbf{p})
  \;=\; C + R_2(\mathbf{p} - C)
  \;=\; R_2(\mathbf{p}) + \underbrace{(C - R_2(C))}_{\Delta}.
\]
The displacement $\Delta = C - R_2(C)$ is a fixed integer vector that depends
only on $C$.  Because the tetrille $v=3$ system is closed under the
\code{tetrille\_translate\_cycle} operation (which works in the coarser $v=6$
units), $\Delta$ must be representable in $v=6$ units.

\subsection{Computing the Translation Constants}

\begin{tcolorbox}[dualbox, title={Translation computation in \code{build\_symmetric\_seed}}]
  The program computes $\Delta$ once at startup by applying the $t=2$
  rotation to a one-vertex cycle containing just $C = v[12]$, then
  computing the displacement back to $C$ in the $v=6$ coordinate system.

  \code{tetrille\_delta\_to\_6(valence, dx, dy, \&m, \&n)} converts an
  integer displacement $(dx, dy)$ (in the coordinate system of a vertex
  with the given valence) into $v=6$ lattice units $(m, n)$.  For
  a $v=3$ vertex: $a = dx - dy$, $b = dx + 2\,dy$; then $m = a/3$,
  $n = b/3$.  The function returns 0 if $dx, dy$ are not exactly
  representable, which would indicate an invalid centre.

  For $v[12] = (3, 1, -1)$: after $t=2$ the point maps to $(3, 0, 1)$,
  giving $dx = 1 - 0 = 1$, $dy = -1 - 1 = -2$.  The formula yields
  $m_6 = +1$, $n_6 = -1$.
\tcblower
\begin{minted}{c}
Coord center = tile->base.v[CENTER_VERTEX];
/* 12 = CENTER_VERTEX */

/* 1-vertex cycle to find where C maps */
Cycle tmp, rtmp;
memset(&tmp, 0, sizeof(tmp));
tmp.n = 1; tmp.v[0] = center;
cycle_transform_lattice(&tmp, &rtmp,
                        tile->lattice, 2);

int dx = center.x - rtmp.v[0].x; /* +1 */
int dy = center.y - rtmp.v[0].y; /* -2 */

tetrille_delta_to_6(center.v, dx, dy,
                    &g_rot_m6, &g_rot_n6);
/* g_rot_m6 = +1, g_rot_n6 = -1 */
\end{minted}
\end{tcolorbox}

\noindent The global pair \code{(g\_rot\_m6, g\_rot\_n6)} is used in every
subsequent call to \code{has\_120\_symmetry}.

% -----------------------------------------------------------------------
\section{Seed Construction}

\subsection{The Three-Hat Pinwheel}

\begin{tcolorbox}[dualbox, title={\code{build\_symmetric\_seed} — overview}]
  Three hat cycles are built at angles $0°$, $120°$, and $240°$:
  \begin{itemize}
    \item \code{rot[0] = tile->base} (orientation $t=0$)
    \item \code{rot[1]}: obtained by \code{cycle\_transform\_lattice(base,\,t=2)}
    \item \code{rot[2]}: obtained by \code{cycle\_transform\_lattice(base,\,t=4)}
  \end{itemize}
  After applying the lattice rotation, $v[12]$ moves to a different position.
  The seed assembly step (below) searches over all $(be, te)$ pairs to find
  the unique attachment of \code{rot[k]} that places its $v[12]$
  back on the shared centre $C$.
\tcblower
\begin{minted}{c}
Cycle rot[SEED_HATS];
rot[0] = tile->base;
cycle_transform_lattice(&tile->base,
    &rot[1], tile->lattice, 2);  /* 120° */
cycle_transform_lattice(&tile->base,
    &rot[2], tile->lattice, 4);  /* 240° */

/* poly[0] = first hat alone */
poly[0].cycle_count = 1;
poly[0].cycles[0]   = rot[0];
hats[0]             = rot[0];
\end{minted}
\end{tcolorbox}

\subsection{Attaching Hat 2 and Hat 3}

\begin{tcolorbox}[dualbox, title={Attachment loop for \code{k = 1, 2}}]
  For each subsequent hat (k=1 and k=2) the program iterates over every
  boundary edge \code{be} of the growing polygon and every edge index
  \code{te} of \code{rot[k]}.  Three conditions must all hold for the
  placement to be accepted:
  \begin{enumerate}
    \item \code{try\_attach\_tile\_poly\_ex} succeeds (the tile can be placed
      without overlap).
    \item The resulting polygon has no holes
      (\code{cycle\_count == 1}).
    \item The placed tile's vertex 12 coincides with the shared centre:
      \code{aligned.v[CENTER\_VERTEX] == center}.
  \end{enumerate}
  The third condition is what enforces the ``three hats all touching at
  $v[12]$'' constraint.  Without it, the loop would accept any non-overlapping
  attachment of the rotated tile.
\tcblower
\begin{minted}{c}
Edge edges[512];
for (int k = 1; k < SEED_HATS; k++) {
    int ok = 0;
    int ec = build_boundary_edges(
                 &poly[k-1], edges);
    for (int be = 0; be < ec && !ok; be++)
    for (int te = 0;
         te < rot[k].n && !ok; te++) {
        Cycle aligned;
        if (!try_attach_tile_poly_ex(
                &poly[k-1], &rot[k],
                tile->lattice, be, te,
                &poly[k], &aligned)) continue;
        if (poly[k].cycle_count > 1) continue;
        /* KEY: v[12] must stay at centre */
        if (!coord_eq(
                aligned.v[CENTER_VERTEX],
                center)) continue;
        hats[k] = aligned;
        ok = 1;
    }
}
\end{minted}
\end{tcolorbox}

\noindent For the canonical hat tile the search succeeds at
$(be=10, te=13)$ for hat 2 and $(be=17, te=13)$ for hat 3, producing
a 30-vertex outer boundary.

\subsection{BComp1 State Initialisation}

Once the three tile cycles are assembled into \code{poly[2]} (a hole-free
polygon) and \code{hats[0..2]}, a \code{BComp1Record} is filled in and
passed to \code{bcomp1\_state\_from\_record} to create the initial growth
state.  Three fields are mandatory for this call to succeed:

\begin{minted}{c}
BComp1Record rec;
memset(&rec, 0, sizeof(rec));
rec.have_boundary = 1;  /* outer polygon */
rec.have_tiles    = 1;  /* tile list      */
rec.have_center   = 1;  /* seed boundary  */
rec.center        = seed_center;     /* = poly[2].cycles[0] */
rec.boundary      = poly[SEED_HATS - 1];
rec.tiles_count   = SEED_HATS;
for (int k = 0; k < SEED_HATS; k++) rec.tiles[k] = hats[k];
return bcomp1_state_from_record(&rec, seed);
\end{minted}

\noindent \textbf{Note:} \code{have\_center} must be set or
\code{bcomp1\_state\_from\_record} returns 0 silently.  The \code{center}
field stores the outer boundary of the 3-hat seed; it is not the abstract
rotation centre $C$ but the \emph{cycle} used by the RL1 engine to track the
boundary of the innermost ``supertile''.

% -----------------------------------------------------------------------
\section{The RL0 / RL1 Growth Engine}

\subsection{BComp1Context}

\begin{tcolorbox}[dualbox, title={\code{BComp1Context} — what is loaded}]
  \code{bcomp1\_context\_init(\&ctx, tile\_path, remembrance\_path, deletions\_path)}
  loads three data sources into a single context object:
  \begin{description}[leftmargin=0pt,style=nextline]
    \item[\code{tile\_path}] The hat tile definition file
      (\file{preferences/focus.tile}), parsed into \code{ctx.tile}
      including all variant cycles.
    \item[\code{remembrance\_path}] The RL0 \emph{remembrance} table
      (\file{data/rl0/remembrance.dat}): the dictionary of valid vertex-arc
      completions at every lattice valence.  Used by \code{attach0\_force\_live\_closure}
      and the BComp1 growth engine to determine which tile placements are
      globally consistent.
    \item[\code{deletions\_path}] The RL0 \emph{deletions} table
      (\file{data/rl0/deletions.dat}): vertex-arc completions that have
      been proved inconsistent and must be excluded.
  \end{description}
  Together the remembrance and deletions tables constitute the \textbf{RL0 rulebook}
  — the first-order local constraint set for valid hat tilings.
\tcblower
\begin{minted}{c}
BComp1Context ctx;
bcomp1_context_init(&ctx,
    "preferences/focus.tile",
    "data/rl0/remembrance.dat",
    "data/rl0/deletions.dat");

/* ctx.tile:  loaded Tile struct    */
/* ctx.map:   RL0 forget/arc map   */
\end{minted}
\end{tcolorbox}

\subsection{Forced Deterministic Closure}

\begin{tcolorbox}[dualbox, title={\code{attach0\_force\_live\_closure}}]
  After the 3-hat pinwheel seed is built, the program attempts a
  \emph{forced ring}: it repeatedly scans every boundary vertex of the
  current polygon and, wherever the RL0 rulebook leaves only one possible
  tile placement, makes that placement deterministically.  This continues
  until no more forced moves remain.

  The function returns 1 (ok) in two situations:
  \begin{itemize}
    \item No forced moves exist from the start (0 steps), or
    \item All forced moves were made and no contradiction was reached.
  \end{itemize}
  It returns 0 only when the RL0 rules produce a contradiction — a
  vertex with no valid completion at all.

  For the $v[12]$ pinwheel, the forced closure completes with
  \code{closure\_steps=0} and \code{unresolved=30}: the 3-hat seed has
  no deterministic forced moves.  Every boundary vertex still has multiple
  valid completions.  This is fundamentally different from the $v[0]$
  pinwheel (where forced closure adds 9 tiles in 6 steps and resolves all
  boundary vertices), and explains why the $v[12]$ seed yields multiple
  symmetric completions at every ring level.
\tcblower
\begin{minted}{c}
Attach0Stats astats;
Attach0ClosureStats cstats;
attach0_stats_init(&astats);
attach0_closure_stats_init(&cstats);

int ok = attach0_force_live_closure(
    &seed.poly,
    &ctx.tile,
    seed.tiles,
    &seed.tile_count,
    &ctx.map,
    FORCE_STEPS,    /* = 1024 */
    &astats,
    &cstats);

/* v[12] result:
   ok=1, closure_steps=0, unresolved=30
   => no forced moves, not a dead end */
\end{minted}
\end{tcolorbox}

\subsection{Ring-by-Ring Growth with BComp1}

\begin{tcolorbox}[dualbox, title={\code{bcomp1\_complete\_state} — growth options}]
  \code{bcomp1\_complete\_state(\&ctx, \&seed, \&center, \&opts, \&result)}
  grows $N$ concentric rings from the seed state.  A ``ring'' is one
  additional layer of RL0-consistent tile placements.  The function
  exhaustively enumerates all completions up to the requested depth and
  collects them into \code{result.records}.

  The three key options are:
  \begin{description}[leftmargin=0pt,style=nextline]
    \item[\code{depth = N}]  Number of rings to add.
    \item[\code{live\_only = 1}]  Discard any completion that the RL0
      rulebook marks as a dead end (a state with no valid infinite
      extension).
    \item[\code{collect\_records = 1}]  Store every surviving completion
      as a \code{BComp1Record} in \code{result.records}.
  \end{description}
  Without \code{live\_only=1} the count at level 3 would be much larger;
  the flag prunes early those clusters that cannot be part of a valid
  hat tiling.
\tcblower
\begin{minted}{c}
BComp1Options opts;
bcomp1_options_default(&opts);
opts.depth           = level; /* --level N */
opts.collect_records = 1;
opts.live_only       = 1;

BComp1Result result;
bcomp1_complete_state(
    &ctx, &seed, &center,
    &opts, &result);

/* result.records.count = total live
   completions at requested depth   */
\end{minted}
\end{tcolorbox}

% -----------------------------------------------------------------------
\section{Symmetry Filtering}

\subsection{The 120° Invariance Test}

\begin{tcolorbox}[dualbox, title={\code{has\_120\_symmetry}}]
  A completion's outer boundary is \textbf{120°-symmetric} if rotating it
  by $120°$ around $C = v[12]$ produces the same cycle (up to a cyclic
  shift of the start vertex).  This is the defining property of the pinwheel
  symmetry group $\mathbb{Z}/3\mathbb{Z}$.

  The test has two steps:
  \begin{enumerate}
    \item Apply the origin-centred $t=2$ lattice rotation to the boundary
      cycle: \code{cycle\_transform\_lattice(c, \&rotated, lattice, 2)}.
    \item Translate the result by \code{(g\_rot\_m6, g\_rot\_n6)} to undo
      the displacement of $C$ under the rotation:
      \code{tetrille\_translate\_cycle(\&rotated, g\_rot\_m6, g\_rot\_n6)}.
  \end{enumerate}
  The two steps together implement $\text{Rot}_{C,120°}$.
  The resulting cycle \code{rotated} is compared to the original \code{c}
  by \code{cycles\_match\_cyclic}, which allows any cyclic shift.
\tcblower
\begin{minted}{c}
static int has_120_symmetry(
    const Poly *boundary, int lattice)
{
    if (boundary->cycle_count != 1)
        return 0;
    const Cycle *c = &boundary->cycles[0];
    Cycle rotated;
    /* step 1: rotate around origin */
    cycle_transform_lattice(
        c, &rotated, lattice, 2);
    /* step 2: translate so C is fixed */
    tetrille_translate_cycle(
        &rotated, g_rot_m6, g_rot_n6);
    /* compare allowing cyclic shift */
    return cycles_match_cyclic(c, &rotated);
}
\end{minted}
\end{tcolorbox}

\subsection{Cyclic Matching}

\begin{tcolorbox}[dualbox, title={\code{cycles\_match\_cyclic}}]
  Two cycles represent the same geometric polygon if they have the same
  vertices in the same order but possibly starting at a different point.
  The function searches all cyclic shifts $s$ of $b$ and checks whether
  any shifted $b$ agrees vertex-for-vertex with $a$.  Both vertex count
  and all coordinate triples $(v, x, y)$ must match exactly.
\tcblower
\begin{minted}{c}
static int cycles_match_cyclic(
    const Cycle *a, const Cycle *b)
{
    if (a->n != b->n) return 0;
    for (int s = 0; s < b->n; s++) {
        if (!coord_eq(b->v[s], a->v[0]))
            continue;       /* wrong start */
        int ok = 1;
        for (int i = 1; i < a->n && ok; i++)
            if (!coord_eq(a->v[i],
                    b->v[(s+i) % b->n]))
                ok = 0;
        if (ok) return 1;
    }
    return 0;
}
\end{minted}
\end{tcolorbox}

\subsection{In-Place Record Filtering}

After \code{bcomp1\_complete\_state} fills \code{result.records}, the
program compacts the array in-place, overwriting non-symmetric entries:

\begin{minted}{c}
size_t total = result.records.count;
size_t sym_count = 0;
for (size_t i = 0; i < total; i++) {
    if (has_120_symmetry(&result.records.items[i].boundary, ctx.tile.lattice))
        result.records.items[sym_count++] = result.records.items[i];
}
result.records.count = sym_count;
\end{minted}

\noindent This avoids any dependency on a non-public \code{push/init/clear}
API for \code{BComp1RecordVec}: the filter works entirely through the public
\code{items[]} and \code{count} fields.

% -----------------------------------------------------------------------
\section{Helper Functions}

\subsection{Collecting Interior (Hidden) Vertices}

\begin{tcolorbox}[dualbox, title={\code{rebuild\_hidden} and \code{collect\_verts}}]
  A \code{BComp1Record} carries a list of \emph{hidden vertices}: lattice
  points that are interior to the cluster (touched by at least one tile
  but not lying on the outer boundary).  The count of hidden vertices
  serves as the RL1 \emph{level} field, a measure of cluster depth.

  \code{collect\_verts} builds the union of all vertex sets from all tile
  cycles (deduplicating by linear scan).  \code{rebuild\_hidden} then
  subtracts the boundary vertices (from \code{build\_boundary\_vertices})
  to obtain only the interior points.  Results are sorted by the
  lexicographic order $(v, x, y)$ for reproducibility.
\tcblower
\begin{minted}{c}
/* Collect all tile vertices */
static int collect_verts(
    const Cycle *tiles, int n, Coord *out)
{
    int c = 0;
    for (int t = 0; t < n; t++)
      for (int i = 0; i < tiles[t].n; i++)
        if (!coord_in(out, c, tiles[t].v[i]))
          out[c++] = tiles[t].v[i];
    return c;
}

/* Hidden = all_verts \ boundary_verts */
static int rebuild_hidden(
    const Poly *p, const Cycle *tiles,
    int n, Coord *hidden)
{
    Coord all[BCOMP1_MAX_COORDS];
    Coord bnd[BCOMP1_MAX_COORDS];
    int ac = collect_verts(tiles, n, all);
    int bc = build_boundary_vertices(p, bnd);
    int hc = 0;
    for (int i = 0; i < ac; i++)
      if (!coord_in(bnd, bc, all[i]))
        hidden[hc++] = all[i];
    qsort(hidden, hc, sizeof(Coord),
          coord_cmp);
    return hc;
}
\end{minted}
\end{tcolorbox}

\subsection{Assembling a BComp1Record from a State}

\code{make\_record(\&state, \&center, \&rec)} constructs the full
\code{BComp1Record} needed to write a structured output file.
It allocates a temporary buffer, calls \code{rebuild\_hidden} to count
and list interior vertices, then populates all mandatory fields:

\begin{minted}{c}
r->have_center = r->have_boundary = r->have_hidden = r->have_tiles = 1;
r->level         = hc;              /* hidden vertex count = depth metric */
r->center        = *center;
r->boundary      = s->poly;
r->hidden_count  = hc;
r->tiles_count   = r->tile_count = s->tile_count;
\end{minted}

% -----------------------------------------------------------------------
\section{SVG Rendering}

\subsection{Coordinate Conversion}

Tetrille coordinates $(v, x, y)$ are mapped to Euclidean screen space using
the basis vectors that the \file{hat.tile} lattice prescribes
($\sqrt{3}$ pre-computed as \code{g\_r3}):

\[
  (x_s,\, y_s) =
  \begin{cases}
    \Bigl(\tfrac{\sqrt{3}}{2}(x+y),\; \tfrac{1}{2}(y-x)\Bigr) & v = 6 \\[4pt]
    \Bigl(\tfrac{\sqrt{3}}{4}(x+y),\; \tfrac{1}{4}(y-x)\Bigr) & v = 4 \\[4pt]
    \Bigl(\tfrac{x + \tfrac{1}{2}y}{\sqrt{3}},\; \tfrac{1}{2}y\Bigr) & v = 3
  \end{cases}
\]

SVG $y$ increases downward, so the draw call negates $y_s$: a tile vertex
at $(x_s, y_s)$ appears at SVG position
$(o_x + \texttt{SC}\cdot x_s,\; o_y - \texttt{SC}\cdot y_s)$
where $(o_x, o_y)$ is the per-cell offset and \code{SC = 24} is the pixel
scale factor.

\subsection{Grid Layout}

A global bounding box is computed over all tile vertices in all output
records, so every grid cell uses the same scale.  The grid dimensions
(rows $\times$ columns) are chosen to minimise the deviation from a $16:9$
aspect ratio:

\begin{minted}{c}
for (int r = 1; r <= n; r++) {
    int c2 = (n + r - 1) / r;
    double sc2 = fabs((double)c2 * cw / ((double)r * ch) - 16.0 / 9.0);
    if (sc2 < best) { best = sc2; rows = r; cols = c2; }
}
\end{minted}

\subsection{Colour Coding}

Tiles are drawn in a fixed layer order (grown tiles first, seed tiles on top,
centre outline last) with the following colour scheme:

\medskip
\begin{tabular}{@{}lll@{}}
\toprule
\textbf{Role} & \textbf{Colour} & \textbf{Hex} \\
\midrule
Seed hat 0 (base orientation)  & {\color{seedorange}\rule{1em}{0.8em}} Amber   & \texttt{\#f39c12} \\
Seed hat 1 (120° rotation)     & {\color{seedred}\rule{1em}{0.8em}} Red     & \texttt{\#e74c3c} \\
Seed hat 2 (240° rotation)     & {\color{seedpurple}\rule{1em}{0.8em}} Purple  & \texttt{\#8e44ad} \\
Grown tiles (ring 1, 2, \ldots) & {\color{grownblue}\rule{1em}{0.8em}} Light blue & \texttt{\#a8d8ea} \\
Seed boundary outline          & {\color{centerred}\rule{1em}{0.8em}} Dark red & \texttt{\#c0392b} \\
\bottomrule
\end{tabular}

\medskip
\noindent The seed boundary outline (the outer boundary of the 3-hat
seed itself) is drawn with no fill and a thick stroke, giving a visual
landmark for the pinwheel centre across all completions.

% -----------------------------------------------------------------------
\section{Main Pipeline}

The \code{main} function executes the following steps in order:

\begin{enumerate}
  \item \textbf{Argument parsing.}  Accepts \code{--level N},
    \code{--tile}, \code{--remembrance}, \code{--deletions},
    \code{--out-dir}, and \code{--svg} flags.
  \item \textbf{Context initialisation.}  \code{bcomp1\_context\_init}
    loads the tile and RL0 rule tables.
  \item \textbf{Seed construction.}  \code{build\_symmetric\_seed}:
    computes rotation constants, builds the three rotated hat cycles,
    searches for the shared-$v[12]$ attachment, assembles the
    \code{BComp1State}.
  \item \textbf{Forced closure.}  \code{attach0\_force\_live\_closure}
    applies all deterministic RL0 moves (0 for the $v[12]$ seed).
  \item \textbf{Supertile record.}  \code{make\_record} + \code{bcomp1\_print\_record}
    writes the seed cluster to \file{data/sym12/supertile.dat}.
  \item \textbf{Level-$N$ growth.}  \code{bcomp1\_complete\_state} with
    \code{depth=N}, \code{live\_only=1}, \code{collect\_records=1}.
  \item \textbf{Symmetry filter.}  In-place compaction retaining only
    records whose boundary satisfies \code{has\_120\_symmetry}.
  \item \textbf{Output.}  Sorted records written to
    \file{data/sym12/sym\_completions\_levelN.dat}; SVG rendered if
    \code{--svg} was given.
\end{enumerate}

% -----------------------------------------------------------------------
\section{Results}

Running \code{./bin/grow\_symmetric2} with the default canonical hat tile
and RL0 rule tables produces:

\medskip
\begin{tabular}{@{}rrrp{8cm}@{}}
\toprule
\textbf{Level} & \textbf{Total live} & \textbf{Symmetric} & \textbf{Remarks} \\
\midrule
1 & 4  & 2  & Two distinct ways to add a first ring with 120° symmetry. \\
2 & 24 & 4  & Four symmetric clusters after two rings. \\
3 & 700 & 10 & Ten symmetric clusters; branching is widening. \\
\bottomrule
\end{tabular}

\medskip
\noindent Diagnostic output to \code{stderr}:

\begin{minted}{text}
Symmetric 3-hat seed (center=v[12]=3(1,-1), rot_m6=1 rot_n6=-1, tiles=3)
grow_symmetric2  center=v[12]  level=3
  forced tiles=3  closure_steps=0  unresolved=30
  completions=700  symmetric=10  ->  data/sym12/sym_completions_level3.dat
\end{minted}

\medskip
\noindent\textbf{Contrast with the $v[0]$ pinwheel.}
\code{grow\_symmetric} (the companion program centred at $v[0]$) gives
1 symmetric completion at levels 1 and 2, and 0 at level 3: the cluster
is a genuine dead end.  The $v[12]$ pinwheel has no such limitation — the
growing count of completions at each level strongly suggests the configuration
can be extended to a valid infinite hat tiling with 3-fold rotational symmetry.

% -----------------------------------------------------------------------
\section{Architecture Summary}

\begin{center}
\begin{tcolorbox}[
  colframe=rulecolor, colback=white, boxrule=0.6pt, arc=3pt,
  width=0.98\linewidth
]
\begin{minted}{text}
main()
  +-- bcomp1_context_init(tile, remembrance, deletions)
  |     loads Tile + RL0 rule tables into BComp1Context
  |
  +-- build_symmetric_seed(ctx.tile, &seed, &center)
  |     +-- center = tile->base.v[12]  = (3,1,-1)
  |     +-- cycle_transform_lattice(tmp, t=2) -> compute g_rot_m6=+1, g_rot_n6=-1
  |     +-- rot[0]=base, rot[1]=t=2(base), rot[2]=t=4(base)
  |     +-- for k=1,2: search (be,te) with aligned.v[12]==center
  |     +-- bcomp1_state_from_record(&rec, &seed)  [need have_center=1]
  |
  +-- attach0_force_live_closure(&seed, ...)
  |     v[12] seed: 0 forced moves, 30 unresolved vertices (open config)
  |
  +-- make_record(&seed, &center, &supertile)
  |     rebuild_hidden: all_verts \ boundary_verts, sorted
  |     -> write data/sym12/supertile.dat
  |
  +-- bcomp1_complete_state(&ctx, &seed, &center, depth=N, live_only=1)
  |     exhaustive ring-by-ring growth guided by RL0 rulebook
  |     result.records.count = total live completions
  |
  +-- symmetry filter (in-place compaction)
  |     for each record: has_120_symmetry(boundary)?
  |       cycle_transform_lattice(t=2) + tetrille_translate_cycle(+1,-1)
  |       cycles_match_cyclic (cyclic shift allowed)
  |     result.records.count = symmetric count
  |
  +-- bcomp1_sort_records + bcomp1_print_record
  |     -> write data/sym12/sym_completions_levelN.dat
  |
  +-- emit_svg (optional)
        bbox over all records -> uniform scale SC=24
        draw grown tiles (#a8d8ea), seed trio (orange/red/purple)
        draw seed boundary outline (#c0392b, no fill)
        grid layout: minimise deviation from 16:9
\end{minted}
\end{tcolorbox}
\end{center}

% -----------------------------------------------------------------------
\section{Usage}

\begin{minted}{bash}
# Build
make grow_symmetric2

# Level 1 (default)
./bin/grow_symmetric2

# Level 3 with SVG output
./bin/grow_symmetric2 --level 3 --svg img/sym12_level3.svg

# Custom output directory and tile file
./bin/grow_symmetric2 --level 2 --out-dir data/my_run \
    --tile preferences/focus.tile \
    --remembrance data/rl0/remembrance.dat \
    --deletions data/rl0/deletions.dat
\end{minted}

\medskip
\noindent Output files written to \code{out-dir} (default \file{data/sym12}):

\begin{itemize}
  \item \file{supertile.dat} — the 3-hat seed cluster as a \code{BComp1Record}.
  \item \file{sym\_completions\_levelN.dat} — all 120°-symmetric live
    completions at the requested ring depth, one record per line.
\end{itemize}

\vspace{4pt}
\noindent\rule{\linewidth}{0.6pt}
\vspace{2pt}

{\small\color{gray}
\file{MyHatTileStuff/grow\_symmetric2.c} — part of PolyformTown.
Compile this document with \texttt{pdflatex -shell-escape grow\_symmetric2.tex}
(requires \texttt{pygments} for \texttt{minted}).
}

\end{document}
